\documentclass[tikz,border=6pt]{standalone}
% Shared visual language for the book's standalone vector figures.
% Keep this file compatible with the TeX Live version used by LyX.
\usepackage{amsmath,amssymb}
\usepackage{array,booktabs}
\usepackage{pgfplots}
\usepackage{pgfplotstable}
\usetikzlibrary{arrows.meta,calc,positioning}
\pgfplotsset{compat=1.17}

\definecolor{BookBlue}{HTML}{0072B2}
\definecolor{BookOrange}{HTML}{D55E00}
\definecolor{BookGreen}{HTML}{009E73}
\definecolor{BookGray}{HTML}{666666}
\definecolor{BookLightGray}{HTML}{D9D9D9}

\pgfplotsset{
  book axis/.style={
    axis line style={black!75,line width=0.55pt},
    tick style={black!75,line width=0.45pt},
    tick label style={font=\small},
    label style={font=\small},
    legend style={
      font=\small,
      draw=none,
      fill=none,
      cells={anchor=west},
    },
    grid style={BookLightGray,line width=0.35pt},
    major grid style={BookLightGray,line width=0.35pt},
  },
}

\begin{document}
\pgfplotstableread[col sep=comma]{../data/rationality_budgets.csv}\budgetdata
\begin{tikzpicture}
  \begin{axis}[
    book axis,
    width=10.5cm,
    height=10.5cm,
    scale only axis,
    axis equal image,
    xmin=0,xmax=100,
    ymin=0,ymax=100,
    xtick={0,10,...,100},
    ytick={0,10,...,100},
    xlabel={$x_1$},
    ylabel={$x_2$},
    grid=major,
  ]
    \pgfplotsinvokeforeach{0,...,49}{
      \pgfplotstablegetelem{#1}{x_start}\of\budgetdata\let\xa\pgfplotsretval
      \pgfplotstablegetelem{#1}{y_start}\of\budgetdata\let\ya\pgfplotsretval
      \pgfplotstablegetelem{#1}{x_end}\of\budgetdata\let\xb\pgfplotsretval
      \pgfplotstablegetelem{#1}{y_end}\of\budgetdata\let\yb\pgfplotsretval
      \edef\drawbudget{\noexpand\addplot[BookGray!48,line width=0.38pt,forget plot]
        coordinates { (\xa,\ya) (\xb,\yb) };}
      \drawbudget
    }
    \addplot[black!65,densely dashed,line width=0.75pt,domain=0:100,samples=2,forget plot] {x};
    \addplot[
      only marks,
      mark=*,
      mark size=1.5pt,
      black,
    ] table[col sep=comma,x=x1,y=x2]{../data/rationality_points.csv};
  \end{axis}
\end{tikzpicture}
\end{document}
